\documentclass[11pt]{article}
\usepackage[a4paper, top=18mm, bottom=18mm, left=20mm, right=20mm]{geometry}
\usepackage[T2A]{fontenc}
\usepackage[utf8]{inputenc}
\usepackage[ukrainian]{babel}
\usepackage{graphicx}
\usepackage[export]{adjustbox}
\usepackage{amsmath}
\usepackage{amssymb}
\graphicspath{{/app/Quiz/Scripts/2023/ProgAll/15BinTree/}{/app/Quiz/Scripts/2021/Mod2021_3/BinTree/}{/app/Quiz/Scripts/2020/Mod2020_4/BinTree/}}
\setlength{\parindent}{0pt}
\setlength{\parskip}{6pt}
\pagestyle{empty}
\begin{document}
%begin
{\large\textbf{Контрольна робота}}

\textbf{Варіант \No\,1}\hfill \textit{ПІБ:}\,\underline{\hspace{60mm}}
\vspace{4mm}

%beginitem
1. %goodpath:Scripts/2018/Mod2018_1/03-good.txt
Відмітити все, що є коректним оператором С++:
( 1 ) return x<=2;

( 2 ) double n.m;

( 3 ) cin>>x+3;

( 4 ) if (1>2) x=10;
\vspace{5mm}
%enditem

%beginitem
2. Що буде виведено за виконання фрагменту коду?

%code
cout << (8 / 8 / 8 / 8);
%code
\vspace{5mm}
%enditem

%beginitem
3. 

a,b,c=1,4,7
def g(b):
global c    a-=5
    b=1
    c=3
    return a+b+c

a,b,c=0,9,5
print(g(a),a,b,c)
\vspace{5mm}
%enditem

%beginitem
4. Що буде виведено за виконання фрагменту коду?

%code
print(3.0 == 3 is 6 != 9)
%code
\vspace{5mm}
%enditem

%beginitem
5. Що буде виведено за виконання фрагменту коду?
%code
a, b, c = 7, 9, 8
def g(a, b, c=6):
    print(a, b, c, end=" ")

g(1, 4, 2)
print(a, b, c)
%code
\vspace{5mm}
%enditem

%beginitem
6. %goodpath:Scripts/2019/Mod2019_1/Lex/01-good.txt
Відмітити все, що є однією правильною лексемою:
( 1 ) 'c'

( 2 ) 0123

( 3 ) 0b23

( 4 ) xy
\vspace{5mm}
%enditem

%beginitem
7. 

print('R', end=' ')
b=['0','1','2','3','4']
b.append('13')
print(b)

\vspace{5mm}
%enditem

%beginitem
8. Що буде виведено за виконання фрагменту коду?

%code
3**-1--1
%code
\vspace{5mm}
%enditem

%beginitem
9. 

\begin{center}
	\adjustimage{max size={0.8\linewidth}{0.8\paperheight}}
	{../15BinTree/trees_exp/35.png}
\end{center}
Указати послідовність міток вершин дерева, зображеного на рис., що утвориться в результаті обходу дерева в оберненому порядку. Мітки записувати через пробіл.\par Алгоритм починає роботу з кореня (на рис. це верхня вершина).
%answer:35 оберненому
\vspace{5mm}
%enditem

%beginitem
10. %goodpath:Scripts/2019/Mod2019_1/Lex/03-good.txt
Відмітити все, що є коректним оператором С++:
( 1 ) return 'a'<'c'

( 2 ) print(sep=":","qwe")

( 3 ) x=*2

( 4 ) float(2)/87
\vspace{5mm}
%enditem

%end
\newpage
\end{document}

